% =============================================================================
% DIM0501 — Boas Práticas de Programação
% Relatório do Projeto da Unidade 1 — AgendaFácil
% Estrutura baseada em Projeto_Unidade1_Modelo_Relatorio.docx
%
% Compilar: pdflatex relatorio.tex   (ou: latexmk -pdf relatorio.tex)
% Limite: o relatório inteiro deve caber em até QUATRO páginas.
% As instruções do modelo original estão em comentários (%) ao longo do
% arquivo e não aparecem no PDF.
% =============================================================================
\documentclass[11pt,a4paper]{article}

\usepackage[utf8]{inputenc}
\usepackage[T1]{fontenc}
\usepackage[brazilian]{babel}
\usepackage{lmodern}
\usepackage[a4paper,margin=2cm]{geometry}
\usepackage{array}
\usepackage{tabularx}
\usepackage{booktabs}
\usepackage{xcolor}
\usepackage{colortbl}
\usepackage{amssymb}
\usepackage{enumitem}
\usepackage{titlesec}
\usepackage{listings}
\usepackage[hidelinks]{hyperref}

% --- Aparência --------------------------------------------------------------
\definecolor{cabecalho}{HTML}{E7EEF7}
\setlength{\parindent}{0pt}
\setlength{\parskip}{4pt}
\renewcommand{\arraystretch}{1.3}
\titleformat{\section}{\large\bfseries}{\thesection.}{0.5em}{}
\titlespacing*{\section}{0pt}{10pt}{4pt}

% Colunas com quebra de linha e texto justificado
% (\hspace{0pt} permite hifenizar a primeira palavra da célula)
\newcolumntype{L}{>{\hspace{0pt}}X}
\newcolumntype{C}[1]{>{\centering\arraybackslash}p{#1}}
\newcommand{\cab}[1]{\cellcolor{cabecalho}\textbf{#1}}

\lstset{
  basicstyle=\ttfamily\small,
  breaklines=true,
  frame=single,
  columns=fullflexible,
  % Permite acentos dentro do bloco de código (saída do programa)
  literate=
    {á}{{\'a}}1 {à}{{\`a}}1 {â}{{\^a}}1 {ã}{{\~a}}1
    {é}{{\'e}}1 {ê}{{\^e}}1 {í}{{\'i}}1
    {ó}{{\'o}}1 {ô}{{\^o}}1 {õ}{{\~o}}1 {ú}{{\'u}}1 {ü}{{\"u}}1 {ç}{{\c{c}}}1
    {Á}{{\'A}}1 {À}{{\`A}}1 {Â}{{\^A}}1 {Ã}{{\~A}}1
    {É}{{\'E}}1 {Ê}{{\^E}}1 {Í}{{\'I}}1
    {Ó}{{\'O}}1 {Ô}{{\^O}}1 {Õ}{{\~O}}1 {Ú}{{\'U}}1 {Ç}{{\c{C}}}1
}

% --- Dados de identificação (preencha aqui) ---------------------------------
\newcommand{\integrantes}{%
  FRANCISCO EDSON DA COSTA FILHO (20240020730) \newline
  JERBESON VINYCIUS SOUSA RODRIGUES (20230025296)%
}
\newcommand{\linguagem}{Python}
\newcommand{\repositorio}{\url{https://github.com/JerbesonViny/boas-praticas-projeto-faculdade}}

\begin{document}

% =============================================================================
% Cabeçalho
% =============================================================================
\begin{center}
  {\small DIM0501 — Boas Práticas de Programação}\\[2pt]
  {\Large\bfseries Relatório do Projeto da Unidade 1 — AgendaFácil}\\[4pt]
  {\small Prof. Marcus Terra \quad·\quad DIMAp/UFRN \quad·\quad
   Semestre 2026.1 \quad·\quad Entrega: 25/10/2026}
\end{center}

\begin{tabularx}{\textwidth}{|>{\raggedright\arraybackslash}p{4.5cm}|L|}
  \hline
  \cab{Integrantes (nome e matrícula)} & \integrantes \\ \hline
  \cab{Linguagem}                       & \linguagem \\ \hline
  \cab{Repositório ou arquivo entregue} & \repositorio \\ \hline
\end{tabularx}

% =============================================================================
\section{Diagnóstico do código-base}
% =============================================================================
% Liste pelo menos dez anomalias encontradas no código-base, com no mínimo uma
% em cada categoria do checklist da Aula 2. Em "onde", aponte a função, a
% variável ou a linha — "está tudo errado" não é localização. Em "por que é um
% problema", escreva o efeito prático sobre quem vai ler ou manter o código,
% não a regra violada. Numere as anomalias: a seção 2 vai se referir a elas
% por número.

\begin{tabularx}{\textwidth}{|C{0.6cm}|>{\raggedright\arraybackslash}p{2.8cm}|L|L|}
  \hline
  \cab{\#} & \cab{Categoria} & \cab{Anomalia e onde está} & \cab{Por que é um problema} \\ \hline
  1  & Programação defensiva
     & Em \texttt{add}, entradas inválidas são ignoradas silenciosamente sob um
       bloco \texttt{try}/\texttt{except: pass} genérico.
     & O chamador acredita que a consulta foi marcada quando na verdade falhou.
       Além disso, oculta bugs internos do próprio código. \\ \hline
  2  & Manutenibilidade
     & Parâmetros com nomes não descritivos em
       \texttt{add(p, pr, d, h, dur)}.
     & Exige leitura interna da função para deduzir os tipos e propósitos de
       cada variável. \\ \hline
  3  & Manutenibilidade
     & Acesso aos dados do agendamento por índices numéricos (ex.:
       \texttt{a[0]}, \texttt{a[1]}, \texttt{a[2]}).
     & Código frágil: alterar a estrutura ou a ordem da lista quebra a aplicação
       silenciosamente em múltiplos pontos. \\ \hline
  4  & Adequação e lógica
     & Ausência de validação para o formato de data (AAAA-MM-DD) nas funções de
       entrada.
     & Permite o cadastro de datas inválidas ou em formatos inconsistentes (ex.:
       \texttt{05/10/2026} ou \texttt{2026-02-30}). \\ \hline
  5  & Adequação e lógica
     & Duração aceita valores arbitrários (não valida o intervalo de 15 a 120 min
       nem múltiplos de 15).
     & Fere as regras de negócio da clínica, permitindo agendamentos com
       durações incoerentes (ex.: 1 min ou 500 min). \\ \hline
  6  & Desempenho e recursos
     & Varredura sequencial direta da lista global \texttt{ag}, duplicada em
       \texttt{add}, \texttt{rem} e \texttt{lst}.
     & Duplicação da lógica de consulta e ausência de uma camada isolada para a
       busca de registros. \\ \hline
  7  & Organização e estética
     & Acoplamento de responsabilidades em \texttt{add} e \texttt{rem}
       (validação, busca, mutação e \texttt{print}).
     & Viola o Princípio da Responsabilidade Única (SRP), dificultando o reuso e
       os testes isolados. \\ \hline
  8  & Programação defensiva
     & Estrutura de \texttt{if} profundamente aninhados dentro de \texttt{add}.
     & Aumenta a complexidade cognitiva e dificulta a manutenção em comparação
       ao uso de cláusulas de guarda. \\ \hline
  9  & Organização e estética
     & Código morto em comentário (\texttt{\# def desmarcar}), TODO antigo e
       cabeçalho com histórico manual.
     & Gera poluição visual. O histórico de alterações pertence ao Git, e não ao
       cabeçalho do arquivo. \\ \hline
  10 & Organização e estética
     & Comentário impreciso em \texttt{lst} (\texttt{\# lista os atendimentos do
       mes}).
     & Comentário incorreto que induz ao erro, pois a função filtra os
       atendimentos por um dia específico. \\ \hline
\end{tabularx}

% =============================================================================
\section{O que foi resolvido e o que ficou como dívida}
% =============================================================================
% Uma linha por anomalia da seção 1, usando o mesmo número. Deixar algo por
% resolver é legítimo — desde que seja uma decisão consciente e justificada, e
% não um esquecimento. Dívida técnica declarada vale mais que dívida técnica
% escondida.
% Preencha a coluna "Situação" com Resolvida ou Dívida assumida.

\begin{tabularx}{\textwidth}{|C{0.6cm}|>{\raggedright\arraybackslash}p{3.2cm}|L|}
  \hline
  \cab{\#} & \cab{Situação} & \cab{Como foi resolvida, ou por que ficou} \\ \hline
  1  & Resolvida
     & Removido o \texttt{try}/\texttt{except: pass}. Cada regra violada lança
       uma exceção específica (\texttt{HorarioForaIntervalo\allowbreak
       PermitidoErro} etc.). \\ \hline
  2  & Resolvida
     & Parâmetros renomeados para \texttt{nome\_paciente},
       \texttt{nome\_doutor}, \texttt{data} etc., seguindo a PEP 8. \\ \hline
  3  & Resolvida
     & Estrutura convertida para a classe \texttt{Agendamento}, com atributos
       nomeados, e o status passou a usar \texttt{AgendamentoStatus}. \\ \hline
  4  & Resolvida
     & Criado \texttt{validar\_data}, que exige estritamente AAAA-MM-DD por
       conversão de ida e volta. \\ \hline
  5  & Resolvida
     & Criado \texttt{validar\_duracao}, que exige de 15 a 120 min, em múltiplos
       de 15 min. \\ \hline
  6  & Dívida assumida
     & Busca mantida por iteração. Como o volume de uma clínica é pequeno, o
       custo é irrisório. Fica centralizada em
       \texttt{\_\_buscar\_agendamento}. \\ \hline
  7  & Resolvida
     & Regras isoladas em \texttt{agenda\_utilitario} e
       \texttt{data\_utilitarios}, aplicadas pela classe \texttt{Agendamento};
       controle de estado na classe \texttt{Agenda}, sem \texttt{print}. \\ \hline
  8  & Resolvida
     & Estrutura alterada para cláusulas de guarda (\emph{guard clauses}), com
       exceções imediatas. \\ \hline
  9  & Resolvida
     & Removidos o TODO, o código comentado e o histórico manual do
       cabeçalho. \\ \hline
  10 & Resolvida
     & Substituído por docstring no padrão Google detalhando a busca por
       data. \\ \hline
\end{tabularx}

% =============================================================================
\section{Decisões de projeto}
% =============================================================================
% Registre de duas a quatro decisões que o código sozinho não explica: por que
% este tipo de dado e não outro, por que lançar exceção em vez de devolver
% valor padrão, por que separar em três funções e não em uma. O que se avalia
% aqui é a justificativa, não a escolha — quase sempre há mais de uma resposta
% defensável.

\begin{tabularx}{\textwidth}{|>{\raggedright\arraybackslash}X|L|L|}
  \hline
  \cab{Decisão tomada} & \cab{Alternativa descartada} & \cab{Por que decidimos assim} \\ \hline
  Validação no construtor do \texttt{Agendamento} (\texttt{se\_validar}).
    & Validar apenas na classe \texttt{Agenda}.
    & Evita a existência de instâncias inconsistentes no sistema. \\ \hline
  Exceções personalizadas, uma por regra.
    & Lançar apenas \texttt{ValueError} genérico ou mensagens de texto.
    & Permite o tratamento granular dos erros pela camada chamadora. \\ \hline
  Validação de data por ida e volta (\texttt{strptime} + \texttt{strftime}).
    & Expressão regular ou \texttt{date.fromisoformat}.
    & Rejeita formatos ambíguos (ex.: \texttt{2026-1-5}) sem a complexidade de
      uma regex. \\ \hline
  Ao remarcar, o status passa a ser \texttt{REMARCADO}.
    & Manter o status \texttt{MARCADO}, como no código-base.
    & A mudança é proposital: o status deve refletir o estado real do
      agendamento. Manter \texttt{MARCADO} após uma remarcação esconderia essa
      informação de quem consulta a agenda. \\ \hline
\end{tabularx}

% =============================================================================
\section{Mudanças de comportamento}
% =============================================================================
% A Etapa 4 muda, de propósito, o que o sistema faz com entradas inválidas.
% Registre cada mudança aqui. O caminho feliz não entra nesta tabela: se ele
% mudou, houve um erro de refatoração.

\begin{tabularx}{\textwidth}{|L|L|>{\raggedright\arraybackslash}X|}
  \hline
  \cab{Entrada ou situação} & \cab{O que acontecia antes} & \cab{O que acontece agora} \\ \hline
  Horário fora do expediente (ex.: 20h)
    & Falhava em silêncio
    & Lança \texttt{HorarioForaIntervalo\allowbreak PermitidoErro} \\ \hline
  Remarcar para horário inválido
    & Remarcava sem validar
    & Lança \texttt{HorarioForaIntervalo\allowbreak PermitidoErro} \\ \hline
  Duração fora de 15 a 120 min
    & Aceitava silenciosamente
    & Lança \texttt{ConsultaIntervalo\allowbreak DuracaoErro} \\ \hline
  Duração não múltipla de 15 min
    & Aceitava
    & Lança \texttt{ConsultaMultiplo\allowbreak DuracaoErro} \\ \hline
  Atendimento que termina depois das 18h (ex.: 17h com 120 min)
    & Aceitava
    & Lança \texttt{ConsultaTermino\allowbreak ForaExpedienteErro} \\ \hline
  Horário ocupado ao marcar ou remarcar
    & Imprimia \texttt{Horario ocupado}
    & Lança \texttt{HorarioReservadoErro} \\ \hline
  Remarcar agendamento inexistente
    & Saía em silêncio
    & Lança \texttt{AgendamentoNao\allowbreak EncontradoErro} \\ \hline
  Data em formato incorreto (ex.: \texttt{2026-1-5})
    & Aceitava como texto solto
    & Lança \texttt{FormatoData\allowbreak InvalidaErro} \\ \hline
  Data inexistente (ex.: \texttt{2026-02-30})
    & Aceitava
    & Lança \texttt{ValueError} \\ \hline
  Data, hora ou duração com tipo errado (ex.: \texttt{None}, \texttt{8.5})
    & Aceitava
    & Lança \texttt{TypeError} \\ \hline
  Erro inesperado interno
    & Engolido por \texttt{except: pass}
    & Propaga a exceção para o chamador \\ \hline
\end{tabularx}

% =============================================================================
\section{Como executar e gerar a documentação}
% =============================================================================

\begin{tabularx}{\textwidth}{|>{\raggedright\arraybackslash}p{5cm}|L|}
  \hline
  \cab{Item} & \cab{Resposta} \\ \hline
  Comando para compilar e executar                 & \texttt{python3 -m src.agenda} \\ \hline
  Comando para gerar a documentação de API         & \texttt{pdoc src -o docs} \\ \hline
  Onde a documentação gerada ficou no repositório  & Diretório \texttt{docs/} \\ \hline
  Saída do programa de demonstração                & (ver abaixo) \\ \hline
\end{tabularx}

\medskip
\textbf{Saída do programa de demonstração (preservada):}
\begin{lstlisting}
Atendimento marcado para Ana
Horario ocupado
Atendimento marcado para Carla
Remarcado
14 Carla Dra. Souza MARCADO
\end{lstlisting}

% =============================================================================
\section{Uso de ferramentas de IA}
% =============================================================================
% O uso é permitido. O professor pode pedir que qualquer integrante explique
% qualquer trecho do código entregue.
% Para marcar: troque $\square$ por $\boxtimes$.

\begin{tabularx}{\textwidth}{|>{\raggedright\arraybackslash}p{5cm}|L|}
  \hline
  \cab{Pergunta} & \cab{Resposta} \\ \hline
  Foram usadas ferramentas de IA neste trabalho?
    & $\square$ Não \qquad $\boxtimes$ Sim \\ \hline
  Em quais etapas, e para quê?
    & Revisão da conformidade com a PEP 8, apoio na estrutura das docstrings
      no padrão Google e formatação do relatório. \\ \hline
\end{tabularx}

% =============================================================================
% Conferência antes de enviar (checklist pessoal — não aparece no PDF)
% =============================================================================
% [ ] O programa de demonstração compila e roda, e imprime a mesma saída do
%     código-base.
% [ ] Toda anomalia da seção 1 aparece na seção 2, resolvida ou declarada como
%     dívida.
% [ ] O README permite que outra pessoa execute o projeto sem perguntar nada.
% [ ] Os textos em itálico deste modelo foram apagados.
%
% "Código que funciona mas não se explica não está pronto."

\end{document}
